%! Author = taj
%! Date = 7/22/25

% Preamble
\documentclass[11pt]{article}

% Packages
\usepackage{amsmath}
\usepackage{graphicx}
\usepackage{float}

\title{8. Pospešeno vrtenje}
\author{Taj Šobot}

% Document
\begin{document}
    \maketitle
    \flushleft

    Pri tem eksperimentu smo preverjali veljavnost 2\. Newtnovega zakona, a v tem primeru za rotacijo.
    Uporabili smo vrteče kolo, na katero je deloval konstantni navor.
    To smo dosegli z padajočo klado.
    Skica eksperimenta:
    \begin{figure}[H]
        \centering
        \includegraphics[width=0.9\textwidth]{img}
        \caption{Skica eksperimenta}
        \label{fig:1}
    \end{figure}

    \section{Meritve}
    Merili smo kako se obhodni čas spreminja skozi pospešeno vrenje.
    Zapisovali smo koliko časa potrebuje vsak zasuk, medtem ko je na obroč deloval konstanten navor.
    Navor je ustvarjala klada, ki je z vrvico bila povezana na sistem.
    Meritve smo opravili za 3 različne mase (navore) in maso obroča:
    \begin{align*}
        m_1 = (50 \pm 1) g \\
        m_2 = (100 \pm 1)g \\
        m_3 = (150 \pm 1)g \\
        m_o = (1160 \pm 1)g
    \end{align*}

    Celoten obroč je imel polmer:
    \begin{align*}
        r_{zunanji} = (25.0 \pm 0.1) cm \\
        r_{notranji} = (25.5 \pm 0.1) cm
    \end{align*}
    Vijak na katerega je delovala radialna sila pa:
    \[
        r_v = (1 cm \pm 0.1) cm
    \]

    Sledi tabela izmerjenih časov:

    \begin{table}[h!]
        \centering
        \begin{tabular}{|c|c|c|c|}
            \hline
            n & $t_1 [s] (m_1)$ & $t_2 [s] (m_2)$ & $t_3 [s] (m_3)$ \\
            \hline
            1	&14.7	&10 	&7.9 \\
            2	&21.1	&14.2	&11.3 \\
            3	&25.9	&17.4	&13.9 \\
            \hline
        \end{tabular}
        \caption{tabela meritev}
        \label{tab:1}
    \end{table}
    Napake:
    \[
        t_{err} = \pm 0.2
    \]

    \section{Rezultati in grafi}
    Definicija navora pravi:
    \[
        \sum \vec{M} = J \vec{\alpha}
    \]
    Iz tega dobimo:
    \[
        F_v r_{v} = J \alpha
    \]
    (r je povprečnavrednost med zunanjim in notranjim)
    \\
    Vstrajnosti moment obroča je:
    \[
        J_o = m r^2 = (0.074 \pm 0.002) kg/m^2
    \]
    Torej:
    \[
        F_v r_{v} / J =  \alpha
    \]
    Tako dobimo kotni pospešek za vsak primer (idealni - izračunan):

    Tabela:
    \begin{table}[h!]
        \centering
        \begin{tabular}{|c|c|}
            \hline
            n & $\alpha [rad/s^2]$\\
            \hline
            1 & 0.068 \\
            2 & 0.135 \\
            3 & 0.202 \\
            \hline
        \end{tabular}
        \caption{tabela meritev}
        \label{tab:2}
    \end{table}

    Ker velja:
    \[
        \phi = (\alpha t^2)/2
    \]
    Oziroma:
    \[
        \alpha = \frac{2 \phi}{t^2}
    \]
    Lahko za vsak izmerjen časa določimo kotni pospešek (ker $\phi = 2\pi$ za vsak zasuk):

    Tabela:
    \begin{table}[h!]
        \centering
        \begin{tabular}{|c|c|c|c|}
            \hline
            $\alpha$ &  $t_{n1} (2\pi)$ & $t_{n2} (4\pi)$& $t_{n3} (6\pi)$\\
            \hline
            $m_1$ & 0.058 & 0.056 & 0.056 \\
            $m_2$ & 0.123 & 0.125 & 0.126 \\
            $m_3$ & 0.201 & 0.197 & 0.195 \\
            \hline
        \end{tabular}
        \caption{tabela kotnih pospeškov}
        \label{tab:5}
    \end{table}

    Povprečni kotni pospešek:
    \begin{table}[h!]
        \centering
        \begin{tabular}{|c|c|c|}
            \hline
            n & $\alpha [rad/s^2]$ (izmerjen) & $\alpha_{err}$\\
            \hline
            1 & 0.057 & 0.004\\
            2 & 0.135 & 0.006\\
            3 & 0.197 & 0.009\\
            \hline
        \end{tabular}
        \caption{tabela povprečji}
        \label{tab:6}
    \end{table}
    Izračun se strinja z meritvami!
    \\

    Izračunajmo še razmerje:
    \[
        t_1^2 : t_2^2 : t_3^2
    \]
    Tabela:
    \begin{table}[H]
        \centering
        \begin{tabular}{|c|c|c|c|}
            \hline
            n & $t_1^2$ & $t_2^2$ & $t_3^2$\\
            \hline
            1 & 216.09 & 100.00 &  62.41 \\
            2 & 445.21 & 201.64 & 127.69 \\
            3 & 670.81 & 302.76 & 193.21 \\
            \hline
        \end{tabular}
        \caption{tabela meritev}
        \label{tab:3}
    \end{table}
    Že iz hitrega pogleda vidimo, da se razmerja ujemajo z 1:2:3 (v stolpcih).

    \section{Izpeljave vprašanja in odgovori}
    Izpeljava za:
    \[
        t_1^2 : t_2^2 : t_3^2 = 1:2:3
    \]
    Velja:
    \begin{align*}
        2 \pi = \frac{\alpha t_1^2}{2} \\
        4 \pi = \frac{\alpha t_2^2}{2} \\
        6 \pi = \frac{\alpha t_3^2}{2} \\
    \end{align*}
    Sledi:
    \begin{align*}
        t_1^2 = \frac{4 \pi }{\alpha}\\
        t_2^2 = \frac{8 \pi }{\alpha}\\
        t_3^2 = \frac{12 \pi }{\alpha}\\
    \end{align*}
    Sledi:
    \begin{align*}
        t_1^2 : t_2^2 : t_3^2 =  \frac{4 \pi }{\alpha}: \frac{8 \pi }{\alpha} : \frac{12 \pi }{\alpha} = 1:2:3
    \end{align*}

    \subsection{Vprašanja}
    \textit{Ocenite, za koliko odstotkov premajhen vztrajnostni moment dobimo, ker v računu
    ne upoštevamo valja in kavljev}
    \\
    Ocenil bi da nekje med 1\% in 2\%.
    
\end{document}